\chapter{Implementation}

\section*{Introduction}

This chapter presents the implementation details of \caimp{}'s key modules,
illustrated with annotated code extracts. We describe the difficulties encountered,
the solutions adopted, and the division of work within the team.

\section{Repository Structure}

The \caimp{} v2 repository uses a monorepo structure, with each service in its
own directory under \texttt{services/}. The 14 database migrations are versioned
under \texttt{infra/postgres/migrations/}, and the shared Python NATS client
library is in \texttt{libs/nats\_client/}.

\section{Work Breakdown}

\begin{table}[H]
\centering
\caption{Team responsibilities}
\label{tab:team}
\small
\begin{tabularx}{\textwidth}{lX}
\toprule
\textbf{Member} & \textbf{Primary modules} \\
\midrule
CHAHIDY Rim  & Telemetry Writer (Go), Alert Engine, React frontend
               (Dashboard, Forecasts, Incidents), CSS design system \\
SAMDI Wiam   & AI Worker, AI Orchestrator, Correlation Engine,
               RAG pipeline, LLM prompts, PDF reports \\
TAÏK Youssef & Admin API, Query API, database (migrations, RLS),
               Change Tracker, CAIMP agent, DevOps (Docker Compose, Makefile, scripts) \\
\bottomrule
\end{tabularx}
\end{table}

\section{Telemetry Writer — Go}

\begin{lstlisting}[language=C, caption={Z-score detector in Go}]
// zScore computes Z-score over the last 100 historical values.
func (d *Detector) zScore(ctx context.Context,
    points []receiver.MetricPoint) []AnomalyEvent {

    var events []AnomalyEvent
    for _, p := range points {
        rows, _ := d.pool.Query(ctx,
            `SELECT value FROM metrics WHERE server_id=$1
             AND metric_name=$2 ORDER BY time DESC LIMIT 100`,
            p.ServerID, p.MetricName)
        var vals []float64
        for rows.Next() {
            var v float64
            if err := rows.Scan(&v); err == nil {
                vals = append(vals, v)
            }
        }
        rows.Close()
        if len(vals) < 30 { continue }  // insufficient history

        mean, std := meanStd(vals)
        if std < 1e-9 { continue }
        z := math.Abs((p.Value - mean) / std)
        if z >= d.cfg.ZScoreThreshold {
            sev := "warning"
            if z >= d.cfg.ZScoreThreshold*1.5 { sev = "critical" }
            events = append(events, AnomalyEvent{
                Detector: "zscore", Severity: sev,
                Value: p.Value, MetricName: p.MetricName,
            })
        }
    }
    return events
}
\end{lstlisting}

\section{Holt-Winters Forecaster — Python/NumPy}

\begin{lstlisting}[language=Python, caption={Holt-Winters double exponential smoothing}]
def _holt_linear(values: np.ndarray, horizon: int,
                 alpha: float = 0.3, beta: float = 0.1
) -> tuple[np.ndarray, float, np.ndarray]:
    """Double exponential smoothing. Returns (forecasts, trend, fitted)."""
    n = len(values)
    level = float(values[0])
    trend = float((values[-1] - values[0]) / max(n - 1, 1))
    fitted = np.empty(n)
    for t in range(n):
        fitted[t] = level + trend
        if t < n - 1:
            prev_level = level
            level = alpha * float(values[t]) + (1 - alpha) * (level + trend)
            trend = beta * (level - prev_level) + (1 - beta) * trend
    # h-step ahead forecast: L_T + h * T_T
    forecasts = np.array([level + (h + 1) * trend for h in range(horizon)])
    return forecasts, trend, fitted
\end{lstlisting}

\section{Change Correlation Scoring}

\begin{lstlisting}[language=Python, caption={Temporal proximity scoring}]
_TYPE_WEIGHT = {
    "deployment": 1.0, "config": 0.9, "package": 0.8,
    "docker_pull": 0.75, "restart": 0.7, "cron": 0.5, "ssh_login": 0.3,
}

def _score_change(change: dict, anomaly_time: datetime) -> float:
    """Score = exp(-0.1 * minutes_before) * type_weight"""
    occurred_at = change["occurred_at"]
    if occurred_at.tzinfo is None:
        occurred_at = occurred_at.replace(tzinfo=timezone.utc)
    minutes_before = (anomaly_time - occurred_at).total_seconds() / 60
    if minutes_before < 0:
        return 0.0
    temporal = math.exp(-0.1 * max(0, minutes_before))
    return temporal * _TYPE_WEIGHT.get(change["change_type"], 0.5)
\end{lstlisting}

\section{Answers-First Dashboard — React}

\begin{lstlisting}[language=C, caption={StatusBanner component in Dashboard.tsx}]
function StatusBanner({ summary }: { summary: DashboardSummary }) {
  const isCritical = summary.overall_status === 'critical'
  const isHealthy  = summary.overall_status === 'healthy'
  const bg = isCritical ? 'var(--danger)'  :
             isHealthy  ? 'var(--success)' : 'var(--warning)'
  return (
    <div style={{ background: faint, border: `1.5px solid ${bg}`,
                  borderRadius: 14, padding: '20px 28px',
                  display: 'flex', alignItems: 'center', gap: 16 }}>
      <div style={{ width: 52, height: 52, borderRadius: '50%',
                    background: bg, display: 'flex',
                    alignItems: 'center', justifyContent: 'center' }}>
        {isCritical ? <AlertTriangle size={28} color="#fff" />
                    : <Check size={28} color="#fff" />}
      </div>
      <div>
        <div style={{ fontSize: 18, fontWeight: 800, color: bg }}>
          {isHealthy ? 'All Systems Healthy' :
           isCritical ? 'Action Required' :
           'Potential Issues Detected'}
        </div>
        <div style={{ fontSize: 13.5, color: 'var(--text-2)' }}>
          {summary.overall_message}
        </div>
      </div>
    </div>
  )
}
\end{lstlisting}

\section{Challenges and Solutions}

\begin{table}[H]
\centering
\caption{Key challenges and solutions}
\label{tab:challenges}
\small
\begin{tabularx}{\textwidth}{lXX}
\toprule
\textbf{Challenge} & \textbf{Impact} & \textbf{Solution} \\
\midrule
OTLP metric name mismatch
  & Agent used \texttt{cpu\_percent}; dashboard expected \texttt{system.cpu.utilization}
  & Lookup table in Query API \\
SSE context corruption in chat
  & Tokens arrived out of order
  & Numbered frames; ordered client buffer \\
pgvector IVFFlat too slow
  & RAG search $>$500ms with 10k vectors
  & Migration to HNSW index, 8x faster \\
Llama 3.1 8B OOM on 4\,GB RAM
  & OOM killer terminates Ollama
  & Test mode: switch to llama3.2:1b \\
Alert duplication flood
  & 50+ identical alerts in 15 minutes
  & Redis NX keys with 15-min TTL \\
\bottomrule
\end{tabularx}
\end{table}

\section*{Conclusion}

This chapter detailed the implementation of \caimp{}'s seven key modules.
The code extracts illustrate the diversity of languages and paradigms mobilised
across the platform. The next chapter evaluates the correctness of these
implementations through systematic testing.
